% =============================================================================
%  SOVEREIGN YIELDS AND THE SPREAD DECOMPOSITION
%
%  INSERT: a new subsection immediately after \subsection{Endogenous Premiums}
%          of 02-model.tex (it continues from eq. bond-decomposition), before
%          \subsection{Government}. It defines eq. hm-yield, which
%          02c-central-bank-tpi.tex now cites rather than redefines.
%
%  Derived from the home banks' bond conditions (eq. bond-pricing), as solved in
%  point_map._bond_demands / _tpi_rule; the yield is the code's
%  delta_b*(1-q)/q (tpi_experiment._sov_bp, reported x4e4 as bp/yr). The
%  additive identities (yield-D), (yield-F) and (spread-decomposition), with
%  the corner term, were checked on 2000 random draws of the quadrature
%  (max error 7e-16). NB the code's own sovereign_spread_legs
%  (output_decomposition.py) is a DIFFERENT, multiplicative split, attributing
%  each leg by removal; no code computes the additive split below yet.
%
%  REVIEW MARKUP (2026-10-07): \yadd / \rdel, macros in the header of
%  monetary_backstop_review.tex.
%
%  NEW SYMBOLS: y^X_t (flow yield), sp_t (spread), R^{b}_{X,t+1} (bond
%    return), \mathrm{el}_t, \mathrm{lp}^X_t, \mathrm{rp}^X_t, \mathrm{cg}^X_t
%    (the four premia, in rate units).
% =============================================================================


\subsection{Sovereign Yields and the Spread}
\label{sec:spread}

Markets quote sovereign risk as yield spreads, not bond prices, and the backstop
of \cref{sec:cb} is stated in spreads too. To close the analysis of bond pricing,
we therefore translate the price decomposition \eqref{eq:bond-decomposition} into
yields and decompose the spread between the two sovereigns in the same way.

\paragraph{The flow yield:}
A unit of the \cite{Hatchondo2009} perpetuity promises the coupon stream
$\delta_b,\,(1-\delta_b)\delta_b,\,(1-\delta_b)^2\delta_b,\dots$ Discounting this
stream at a constant rate $y$ gives
$q=\sum_{k\ge0}\delta_b(1-\delta_b)^k/(1+y)^{k+1}=\delta_b/(\delta_b+y)$, so the
flow yield of a bond priced at $q$ is
\begin{equation}
y(q)\;=\;\delta_b\,\frac{1-q}{q},
\qquad
y^{X}_t\equiv y\bigl(q^{X}_t\bigr),
\label{eq:hm-yield}
\end{equation}
and the sovereign spread is the gap between the two yields:
\begin{equation}
sp_t\;\equiv\;y^{D}_t-y^{F}_t .
\label{eq:spread-def}
\end{equation}
All rates are quarterly, and we report them annualised in basis points.
Differentiating \eqref{eq:hm-yield} gives the bond's modified duration,
$-\partial\log q/\partial y=1/(\delta_b+y)$, about four years at the calibrated
$\delta_b$. Since $\delta_b+y^{X}_t=\delta_b/q^{X}_t$, the spread and the two
prices are linked exactly by
\begin{equation}
sp_t\;=\;\bigl(\delta_b+y^{F}_t\bigr)\left(\frac{q^{F}_t}{q^{D}_t}-1\right),
\qquad
\log\frac{\delta_b+y^{D}_t}{\delta_b+y^{F}_t}\;=\;\log q^{F}_t-\log q^{D}_t .
\label{eq:spread-prices}
\end{equation}
The spread is the relative price gap divided by the duration of the safe bond:
a $D$ bond priced 1\% below the $F$ bond trades at roughly 25 basis points a year
of spread. The log form shows that any multiplicative decomposition of the two
prices carries over, leg by leg, to the gross flow yields $\delta_b+y$. The
yields themselves, however, are not linear in the price, so the additive
decomposition \eqref{eq:bond-decomposition} does not carry over directly. We
derive its counterpart in rate units from the return on the bond.

\paragraph{The required return:}
Applied to the domestic sovereign, the bank's pricing condition sets the
realised gross return $R^{b}_{D,t+1}=\Xi^{D}_{t+1}/q^{D}_t$ against the deposit
rate. Dividing \eqref{eq:euler-rearranged} by
$q^{D}_t\,\E_t\bigl[\Omega^{(d')}_{D,t,t+1}\bigr]$ and applying the covariance
identity, the bond's expected return is
\begin{equation}
\E_t\bigl[R^{b}_{D,t+1}\bigr]
\;=\;R_{D,t}
\;+\;\frac{\lambda_D\,\mu_{D,t}}{\E_t\bigl[\Omega^{(d')}_{D,t,t+1}\bigr]}
\;-\;\frac{\mathrm{Cov}_t\bigl(\Omega^{(d')}_{D,t,t+1},\,\Xi^{D}_{t+1}\bigr)}
         {q^{D}_t\,\E_t\bigl[\Omega^{(d')}_{D,t,t+1}\bigr]} .
\label{eq:required-return}
\end{equation}
Equivalently, \eqref{eq:required-return} is \eqref{eq:bond-decomposition}
multiplied through by $R_{D,t}/q^{D}_t$. The bond's expected return has to cover
the bank's funding cost, a liquidity premium and a risk premium. The liquidity
premium is exactly the credit spread of \eqref{eq:rwc}. Under the single
divertibility $\lambda_D$ of \eqref{eq:IC}, every asset on the bank's book earns
the same constraint premium as the working-capital loan, and the sovereign bond
is no exception.

\paragraph{From the return to the yield:}
The Hatchondo--Martinez payoff splits the return into a coupon yield and a
capital gain, both scaled by the survival factor. Since
$\delta_b/q^{D}_t=\delta_b+y^{D}_t$,
\begin{equation}
R^{b}_{D,t+1}
\;=\;h_{t+1}\left[\,1+y^{D}_t
+(1-\delta_b)\,\frac{q^{D}_{t+1}-q^{D}_t}{q^{D}_t}\right].
\label{eq:hm-return}
\end{equation}
Under \eqref{eq:expectation} the survival factor has mean
$\E_t[h_{t+1}]=1-\pi^{d}_t(1-\varrho_D)$. Taking expectations of
\eqref{eq:hm-return}, substituting into \eqref{eq:required-return} and solving for
the yield gives the decomposition of the $D$ yield:
\begin{equation}
y^{D}_t
\;=\;\underbrace{r_{D,t}}_{\text{deposit rate}}
\;+\;\underbrace{\pi^{d}_t\,(1-\varrho_D)\,\bigl(1+y^{D}_t\bigr)}_{\mathrm{el}_t:\ \text{expected default loss}}
\;+\;\underbrace{\frac{\lambda_D\,\mu_{D,t}}{\E_t[\Omega_D]}}_{\mathrm{lp}^{D}_t:\ \text{liquidity premium}}
\;\underbrace{-\;\frac{\mathrm{Cov}_t\bigl(\Omega^{(d')}_{D,t,t+1},\,\Xi^{D}_{t+1}\bigr)}
         {q^{D}_t\,\E_t[\Omega_D]}}_{\mathrm{rp}^{D}_t:\ \text{risk premium}\;(>0)}
\;\underbrace{-\;(1-\delta_b)\,\frac{\E_t\bigl[h_{t+1}\bigl(q^{D}_{t+1}-q^{D}_t\bigr)\bigr]}{q^{D}_t}}_{-\mathrm{cg}^{D}_t:\ \text{expected capital gain}} .
\label{eq:yield-decomposition-D}
\end{equation}
The identity is exact, and, like \eqref{eq:bond-decomposition}, it holds wherever
the bank's pricing condition does. As $q^{D}_t$ did there, $y^{D}_t$ appears on
both sides. Each of its five terms has a distinct economic content.
\begin{enumerate}
\item The \emph{deposit rate} is the bank's marginal funding cost, the
  benchmark against which every asset on its book is priced.
\item The \emph{expected default loss} is the probability of default times the
  loss given default, applied to the value of the claim at an unchanged price.
  It is actuarial: it contains no discount factor, and it is the only premium
  over the deposit rate that a risk-neutral, unconstrained investor would demand
  in a stationary state.
\item The \emph{liquidity premium} is the shadow cost of the incentive
  constraint per unit of assets. It is the yield counterpart of the
  liquidity discount in \eqref{eq:bond-decomposition}, and it is non-zero only
  when the constraint binds today, $\mu_{D,t}>0$.
\item The \emph{risk premium} is the yield counterpart of the covariance term in
  \eqref{eq:bond-decomposition}. The bond pays least in the default branch,
  where the banker values wealth most, $\Omega^{(d=1)}>\Omega^{(d=0)}$, so the
  covariance is negative and the premium positive. It operates even when
  today's constraint is slack.
\item The \emph{expected capital gain} enters because the flow yield counts only
  the coupon. When the price is expected to move, part of the required return is
  delivered as a capital gain or loss, and the flow yield adjusts by the
  opposite amount. After a rise in $s_t$ the price is depressed and expected to
  recover as $s_t$ reverts to $\bar s$ in \eqref{eq:s-process}, so this term
  pulls the flow yield below the required return.
\end{enumerate}

\paragraph{The safe yield:}
The $F$ bond is priced by the $F$ bank, \eqref{eq:bond-pricing}, and never
defaults. The same steps with $h\equiv1$ give
\begin{equation}
y^{F}_t
\;=\;r_{F,t}
\;+\;\underbrace{\frac{\lambda_F\,\mu_{F,t}}{\E_t[\Omega_F]}}_{\mathrm{lp}^{F}_t}
\;\underbrace{-\;\frac{\mathrm{Cov}_t\bigl(\Omega^{(d')}_{F,t,t+1},\,\Xi^{F}_{t+1}\bigr)}
         {q^{F}_t\,\E_t[\Omega_F]}}_{\mathrm{rp}^{F}_t}
\;\underbrace{-\;(1-\delta_b)\,\frac{\E_t\bigl[q^{F}_{t+1}-q^{F}_t\bigr]}{q^{F}_t}}_{-\mathrm{cg}^{F}_t} .
\label{eq:yield-decomposition-F}
\end{equation}
The $F$ bond carries no default loss, but it need not carry a zero risk premium.
The $F$ bank holds $D$ bonds too, so its kernel is also higher in the default
branch, $\Omega^{(d=1)}_{F}>\Omega^{(d=0)}_{F}$. If the $F$ bond appreciates in
that branch, as it \rdel{does}\yadd{would} in a flight to safety, the covariance is positive and
$\mathrm{rp}^{F}_t<0$: the safe bond is a hedge, and its yield falls below the
$F$ deposit rate plus the liquidity premium.

\paragraph{The spread decomposition:}
Subtracting \eqref{eq:yield-decomposition-F} from
\eqref{eq:yield-decomposition-D},
\begin{equation}
sp_t
\;=\;\underbrace{\bigl(r_{D,t}-r_{F,t}\bigr)}_{\text{deposit-rate gap}}
\;+\;\underbrace{\mathrm{el}_t}_{\text{expected loss}}
\;+\;\underbrace{\bigl(r^{wc}_{D,t}-r_{D,t}\bigr)-\bigl(r^{wc}_{F,t}-r_{F,t}\bigr)}_{\text{liquidity: the credit-spread gap}}
\;+\;\underbrace{\bigl(\mathrm{rp}^{D}_t-\mathrm{rp}^{F}_t\bigr)}_{\text{risk-premium gap}}
\;-\;\underbrace{\bigl(\mathrm{cg}^{D}_t-\mathrm{cg}^{F}_t\bigr)}_{\text{capital-gain gap}},
\label{eq:spread-decomposition}
\end{equation}
where the liquidity term uses \eqref{eq:rwc},
$\mathrm{lp}^{X}_t=r^{wc}_{X,t}-r_{X,t}$. The decomposition highlights that the
sovereign spread is more than the price of expected default, and it reads term
by term.

\emph{The deposit-rate gap.} Each bond is a claim on its issuer's good, so the
two yields are own-good real yields. By the parity condition \eqref{eq:uip},
$r_{D,t}-r_{F,t}=(1+r_{F,t})\bigl(\E_t[p_{t+1}]/p_t-1\bigr)
-\kappa_{\rm nfa}V^{\rm flow}_{D,t}/\bar Y_D$: the expected change in the terms of
trade plus the external premium. This is the real counterpart of the inflation
differential in a nominal union. It is zero at the symmetric steady state and
opens only off it.

\emph{The expected loss.} In a stationary state with risk-neutral, unconstrained
investors and a common deposit rate, every other term vanishes, and
\eqref{eq:spread-decomposition} collapses to the textbook credit triangle,
$sp_t=\pi^{d}_t(1-\varrho_D)(1+y^{D}_t)$. Everything else in the spread comes from
intermediation or from the dynamics of prices.

\emph{The liquidity term.} The liquidity component of the sovereign spread is
\emph{exactly} the gap between the two countries' credit spreads. This is a
direct implication of the single $\lambda$: when sovereign losses tighten the $D$
constraint, the same change in $\lambda_D\mu_{D,t}/\E_t[\Omega_D]$ raises both
the lending rate firms pay in \eqref{eq:labour-demand} and the yield the $D$
sovereign pays on its debt. Through this term the sovereign spread and the cost of
credit to firms move together by construction, which is one face of the
pass-through of \cite{bocola2016pass}, written in yields.

\emph{The risk-premium gap.} Here the anticipation of default is priced through
the bankers' kernels rather than through today's constraint. It widens the
spread through both legs: the $D$ bond is a poor asset in the default branch,
and the $F$ bond may be a good one.

\emph{The capital-gain gap.} This is the term through which duration and the
persistence of the risk factor act\yadd{ most directly; they also enter the risk premium, through the resale value $q^{D}_{t+1}$ inside $\Xi^{D}_{t+1}$}. A long bond whose price is expected to
recover delivers part of its return as a capital gain, so its flow yield
understates its required return. How much depends on how far the price has
fallen and how fast \eqref{eq:s-process} brings it back.

The five terms form an accounting identity, not five independent channels. In
equilibrium they move jointly, and none of them can be switched off while the
others are held fixed. \yadd{The risk premium and the capital gain show why. Holding next period's prices fixed, a one-period rise in $\mathrm{rp}^{D}_t$ lowers $q^{D}_t$ and raises $\mathrm{cg}^{D}_t$ almost one for one, so the flow yield rises by only $(\delta_b+y^{D}_t)/(1+y^{D}_t)$ times the premium, about 6\% of it at the calibrated $\delta_b$. A persistent premium moves the yield fully, because it lowers $q^{D}_{t+1}$ as well.} What \eqref{eq:spread-decomposition} delivers is the
mechanism behind each basis point of the spread, not a counterfactual.

\paragraph{The backstop wedge:}
Under the backstop of \cref{sec:cb} the $D$ bank's condition can hold with
inequality, $\mathcal{F}_{D,t}<0$ in \eqref{eq:bank-bond-kkt}, when the floor
binds and the bank sells its whole book. The same derivation then adds one term
to \eqref{eq:yield-decomposition-D},
\begin{equation}
y^{D}_t\;=\;\Bigl[r_{D,t}+\mathrm{el}_t+\mathrm{lp}^{D}_t+\mathrm{rp}^{D}_t-\mathrm{cg}^{D}_t\Bigr]
\;+\;\frac{\mathcal{F}_{D,t}}{q^{D}_t\,\E_t[\Omega_D]},
\qquad
\frac{\mathcal{F}_{D,t}}{q^{D}_t\,\E_t[\Omega_D]}\;\le\;0 ,
\label{eq:yield-wedge}
\end{equation}
by which the central bank holds the yield below the return the bank would require
to hold the bond. Wherever the floor binds, the spread equals the cap,
$sp_t=\bar c$. If the bank still holds part of its book, $\mathcal{F}_{D,t}=0$ and
\eqref{eq:spread-decomposition} shows which premia the purchase compressed to
meet the cap. At the corner, the wedge in \eqref{eq:yield-wedge} absorbs the
remainder.
